\documentclass[answers]{exam}
\usepackage{../../../mypackages}
\usepackage{../../../macros}


\tcbuselibrary{theorems,skins,hooks}

\definecolor{myindbg}{HTML}{DFF5EC}
\definecolor{myindfr}{HTML}{1F6F4A}

\newtcolorbox{Indication}{
  enhanced,
  colback=myindbg,
  colframe=myindfr,
  arc=6mm,
  rounded corners,
  boxrule=0.6mm,
  left=3mm,
  right=3mm,
  top=2mm,
  bottom=2mm
}

\title{Interrogation N°5}
\author{N. Bancel}
\date{3 Février 2026}

% Mise en forme des solutions en bleu
\SolutionEmphasis{\color{blue}}
\renewcommand{\solutiontitle}{\noindent}

\begin{document}

\dsheader{Collège Lycée Suger}{Mathématiques}{Année 2025-2026}{Classe de 6ème}
{\let\newpage\relax\maketitle}
\consignesrendusujet{45 minutes}{0.5}

\section*{Exercice 1 - Cours - Médiatrice et bissectrice (4 points)}

\begin{questions}
  \question[2] \textbf{Médiatrice} Donner la définition précise de la médiatrice d’un segment. Un vocabulaire mathématique est attendu.
  \begin{solution}
          \begin{Indication}

Ce qui est attendu : 
\begin{compactitem}
\item Penser à la perpendicularité
\item Penser au milieu du segment
\item Eviter les tournures "la droite qui forme un angle droit". "La droite fait 90° d'angle". Forme un angle droit avec quoi ?
\item Des points symétriques correctement placés : avec le quadrillage, il n'y a pas vraiment d'excuse pour un mauvais placement et des traits pas droits.
\end{compactitem}

      \end{Indication}
      La médiatrice d’un segment est \textbf{la droite qui passe par le milieu de ce segment et qui est perpendiculaire à ce segment}.
      \end{solution}
  \question[2] \textbf{Bissectrice} Donner la définition précise de la bissectrice d’un angle. Un vocabulaire mathématique est attendu.
      \begin{solution}
                  \begin{Indication}

Ce qui est attendu : 
\begin{compactitem}
\item On ne parle pas du milieu d'un angle  
\item On ne dit pas que l'angle est coupé en deux "parties égales" (trop vague). On dit qu'il est partagé en deux angles de même mesure.
\item On évite de dire que les angles sont égaux. Au fond, c'est leur mesure qui est égale.
\end{compactitem}

      \end{Indication}   
      La bissectrice d’un angle est \textbf{la demi-droite qui partage cet angle en deux angles de même mesure}.
    \end{solution}
  \end{questions}

\section*{Exercice 2 - Symétrie axiale - Exercice basique (7 points)}

\begin{questions}
  \question[2] Sur la figure ci-dessous, tracer le symétrique de la figure dessinée à gauche par rapport à la droite $(d)$. Il n'est pas nécessaire ici de mettre de légendes (les tracés sont à faire sur le sujet).
\end{questions}

\begin{solution}
      \begin{Indication}

Ce qui est attendu : 
\begin{compactitem}
\item Une figure proprement tracée
\item Des points symétriques correctement placés : avec le quadrillage, il n'y a pas vraiment d'excuse pour un mauvais placement et des traits pas droits.
\end{compactitem}

      \end{Indication}

      \fig{0.8}{corr_01.jpg}{}

      \end{solution}

\begin{questions}
  \question[5] Sur la figure ci-dessous, tracer le symétrique du cercle $C$ et de rayon $O$ par rapport à la droite $(d1)$, et tracer le symétrique de la droite $(d2)$ par rapport à la droite $(d1)$. \textrr{Vous noterez toutes les légendes nécessaires} (les tracés sont à faire sur le sujet).
\fig{1}{00.jpg}{}
\end{questions}

\begin{solution}
      \begin{Indication}

Ce qui est attendu : 
\begin{compactitem}
\item Que les angles droits soient faits à l'équerre
\item Que les traits utilisés pour arriver aux figures soient moins visibles que les figures elles-mêmes (vous pouvez appuyer moins fort ou faire des pointillés)
\item Que le cercle soit bien tracé au compas
\item Que toutes les légendes demandées soient présentes : quand on fait le symétrique d'un point, on nomme le point d'origine et on nomme son symétrique. Et on fait les codages de longueur et d'angle qui nous ont permis de tracer le symétrique.
\item \textrr{Il n'y a pas de codage de longueur à faire sur une droite. Une droite n'a pas de longueur}
\item \textrr{Une droite a besoin d'un 2ème point pour être tracée. Il fallait choisir un 2ème point au hasard sur la 1ère droite, le nommer, et construire son symétrique}
\end{compactitem}

      \end{Indication}
      \fig{1}{corr_02.jpg}{}


\end{solution}

\section*{Exercice 4 - Axes de symétrie (5 points)}




\begin{questions}
  \question[5] Sans justifier, tracer les axes de symétrie des figures ci-dessous (les tracés sont à faire sur le sujet).

\end{questions}

\begin{solution}
  
\begin{Indication}

Ce qui est attendu : 
\begin{compactitem}
\item Idéalement : tracer des droites, pas des segments qui s'arrêtent pile aux extremités des panneaux.
\item Tracer les axes à la règle, même si c'est approximatif. Et propremetn
\end{compactitem}

\end{Indication}

      \fig{0.8}{corr_03.jpg}{}
\end{solution}

\section*{Exercice 5 - Déductions (4 points)}

\begin{questions}
  \question[4] Le polygône $DEF$ a pour symétrique par rapport à la droite $(d)$ le polygône $D'E'F'$. $I'$ est le symétrique du point $I$ par rapport à la droite $(d)$. $J \in (DE)$ et $J'$ est le symétrique du point $J$ par rapport à la droite $(d)$. On considérera que les longueurs affichées correspondent à des $cm$, et que $de$ correspond à la longueur du segment $[DE]$,
  que $ef$ correspond à la longueur du segment $[EF]$ 
  et que $df$ correspond à la longueur du segment $[DF]$.
  \fig{0.8}{03.jpg}{}
  \begin{Indication}

Ce qui est attendu : 
\begin{compactitem}
\item Une structure claire de réponse avec le classique "On sait que (hypothèse) Or (propriété) Donc (conclusion)"
\item Perte de la moitié des points si la propriété de conservation soit des angles, soit des longueurs, soit des milieux, soit des alignements n'est pas citée
\item Perte de la moitié des points si dans le "on sait que", on ne cite pas les points importants pour la résolution de la question, et si on ne mentionne pas explicitement qu'ils sont symétriques.
\item Faire une conclusion : on attend des valeurs.
\item Ce ne sont pas des longueurs ou des mesures qui sont symétriques. Ce sont des points, ou des segments, ou des angles, ou des droites etc (bref, des objets, pas des valeurs)
\end{compactitem}

\end{Indication}
  \begin{parts}
    \part[1] Que peut-on dire de l'angle $\widehat{E'D'F'}$ ? Justifier.
    \begin{solution}
On sait que le polygône $DEF$ a pour symétrique par rapport à la droite $(d)$ le polygône $D'E'F'$, donc les points $D,E,F$ ont pour images $D',E',F'$.
L'angle $\widehat{E'D'F'}$ est donc l'image de l'angle $\widehat{EDF}$ par la symétrie axiale d'axe $(d)$.
Or la symétrie axiale conserve les angles.
Donc $\widehat{E'D'F'} = \widehat{EDF} = 60{,}81^\circ$.
\end{solution}
    \part[1] Que peut-on dire de la longueur $D'F'$ ? Justifier.
        \begin{solution}
On sait que $[D'F']$ est l'image du segment $[DF]$ par la symétrie axiale d'axe $(d)$.
Or la symétrie axiale conserve les longueurs.
Donc $D'F' = DF = 9{,}28\ \text{cm}$
\end{solution}
    \part[1] Que peut-on dire du point $I'$ par rapport au segment $[E'F']$ ? Justifier.
            \begin{solution}
On sait que $I$ est le milieu du segment $[EF]$, d'après le codage de la figure. $I', E', D'$ sont les symétriques respectifs de $I, E, F$ par rapport à la droite $(d)$.
Or la symétrie axiale conserve les milieux
Donc $I'$ est le milieu du segment $[E'F']$.
\end{solution}
    \part[1] Que peut-on dire de l'angle $\widehat{E'D'J'}$ ? Justifier.
    \begin{solution}
On sait que $J \in (DE)$ donc les points J, D et E sont alignés. On sait aussi que $J', D', E'$ sont les symétriques respectifs de $J, D, E$ par rapport à la droite $(d)$.
Or la symétrie axiale conserve les alignements, donc les points $D',E',J'$ sont alignés.
On peut donc dire que l'angle $\widehat{E'D'J'}$ est un angle plat, c'est-à-dire que $\widehat{E'D'J'} = 180^\circ$.
\end{solution}
  \end{parts}

\end{questions}

\end{document}
